% Options for packages loaded elsewhere
% Options for packages loaded elsewhere
\PassOptionsToPackage{unicode}{hyperref}
\PassOptionsToPackage{hyphens}{url}
\PassOptionsToPackage{dvipsnames,svgnames,x11names}{xcolor}
%
\documentclass[
  a4paper,
]{article}
\usepackage{xcolor}
\usepackage{amsmath,amssymb}
\setcounter{secnumdepth}{-\maxdimen} % remove section numbering
\usepackage{iftex}
\ifPDFTeX
  \usepackage[T1]{fontenc}
  \usepackage[utf8]{inputenc}
  \usepackage{textcomp} % provide euro and other symbols
\else % if luatex or xetex
  \usepackage{unicode-math} % this also loads fontspec
  \defaultfontfeatures{Scale=MatchLowercase}
  \defaultfontfeatures[\rmfamily]{Ligatures=TeX,Scale=1}
\fi
\usepackage{lmodern}
\ifPDFTeX\else
  % xetex/luatex font selection
\fi
% Use upquote if available, for straight quotes in verbatim environments
\IfFileExists{upquote.sty}{\usepackage{upquote}}{}
\IfFileExists{microtype.sty}{% use microtype if available
  \usepackage[]{microtype}
  \UseMicrotypeSet[protrusion]{basicmath} % disable protrusion for tt fonts
}{}
\makeatletter
\@ifundefined{KOMAClassName}{% if non-KOMA class
  \IfFileExists{parskip.sty}{%
    \usepackage{parskip}
  }{% else
    \setlength{\parindent}{0pt}
    \setlength{\parskip}{6pt plus 2pt minus 1pt}}
}{% if KOMA class
  \KOMAoptions{parskip=half}}
\makeatother
% Make \paragraph and \subparagraph free-standing
\makeatletter
\ifx\paragraph\undefined\else
  \let\oldparagraph\paragraph
  \renewcommand{\paragraph}{
    \@ifstar
      \xxxParagraphStar
      \xxxParagraphNoStar
  }
  \newcommand{\xxxParagraphStar}[1]{\oldparagraph*{#1}\mbox{}}
  \newcommand{\xxxParagraphNoStar}[1]{\oldparagraph{#1}\mbox{}}
\fi
\ifx\subparagraph\undefined\else
  \let\oldsubparagraph\subparagraph
  \renewcommand{\subparagraph}{
    \@ifstar
      \xxxSubParagraphStar
      \xxxSubParagraphNoStar
  }
  \newcommand{\xxxSubParagraphStar}[1]{\oldsubparagraph*{#1}\mbox{}}
  \newcommand{\xxxSubParagraphNoStar}[1]{\oldsubparagraph{#1}\mbox{}}
\fi
\makeatother

\usepackage{color}
\usepackage{fancyvrb}
\newcommand{\VerbBar}{|}
\newcommand{\VERB}{\Verb[commandchars=\\\{\}]}
\DefineVerbatimEnvironment{Highlighting}{Verbatim}{commandchars=\\\{\}}
% Add ',fontsize=\small' for more characters per line
\usepackage{framed}
\definecolor{shadecolor}{RGB}{241,243,245}
\newenvironment{Shaded}{\begin{snugshade}}{\end{snugshade}}
\newcommand{\AlertTok}[1]{\textcolor[rgb]{0.68,0.00,0.00}{#1}}
\newcommand{\AnnotationTok}[1]{\textcolor[rgb]{0.37,0.37,0.37}{#1}}
\newcommand{\AttributeTok}[1]{\textcolor[rgb]{0.40,0.45,0.13}{#1}}
\newcommand{\BaseNTok}[1]{\textcolor[rgb]{0.68,0.00,0.00}{#1}}
\newcommand{\BuiltInTok}[1]{\textcolor[rgb]{0.00,0.23,0.31}{#1}}
\newcommand{\CharTok}[1]{\textcolor[rgb]{0.13,0.47,0.30}{#1}}
\newcommand{\CommentTok}[1]{\textcolor[rgb]{0.37,0.37,0.37}{#1}}
\newcommand{\CommentVarTok}[1]{\textcolor[rgb]{0.37,0.37,0.37}{\textit{#1}}}
\newcommand{\ConstantTok}[1]{\textcolor[rgb]{0.56,0.35,0.01}{#1}}
\newcommand{\ControlFlowTok}[1]{\textcolor[rgb]{0.00,0.23,0.31}{\textbf{#1}}}
\newcommand{\DataTypeTok}[1]{\textcolor[rgb]{0.68,0.00,0.00}{#1}}
\newcommand{\DecValTok}[1]{\textcolor[rgb]{0.68,0.00,0.00}{#1}}
\newcommand{\DocumentationTok}[1]{\textcolor[rgb]{0.37,0.37,0.37}{\textit{#1}}}
\newcommand{\ErrorTok}[1]{\textcolor[rgb]{0.68,0.00,0.00}{#1}}
\newcommand{\ExtensionTok}[1]{\textcolor[rgb]{0.00,0.23,0.31}{#1}}
\newcommand{\FloatTok}[1]{\textcolor[rgb]{0.68,0.00,0.00}{#1}}
\newcommand{\FunctionTok}[1]{\textcolor[rgb]{0.28,0.35,0.67}{#1}}
\newcommand{\ImportTok}[1]{\textcolor[rgb]{0.00,0.46,0.62}{#1}}
\newcommand{\InformationTok}[1]{\textcolor[rgb]{0.37,0.37,0.37}{#1}}
\newcommand{\KeywordTok}[1]{\textcolor[rgb]{0.00,0.23,0.31}{\textbf{#1}}}
\newcommand{\NormalTok}[1]{\textcolor[rgb]{0.00,0.23,0.31}{#1}}
\newcommand{\OperatorTok}[1]{\textcolor[rgb]{0.37,0.37,0.37}{#1}}
\newcommand{\OtherTok}[1]{\textcolor[rgb]{0.00,0.23,0.31}{#1}}
\newcommand{\PreprocessorTok}[1]{\textcolor[rgb]{0.68,0.00,0.00}{#1}}
\newcommand{\RegionMarkerTok}[1]{\textcolor[rgb]{0.00,0.23,0.31}{#1}}
\newcommand{\SpecialCharTok}[1]{\textcolor[rgb]{0.37,0.37,0.37}{#1}}
\newcommand{\SpecialStringTok}[1]{\textcolor[rgb]{0.13,0.47,0.30}{#1}}
\newcommand{\StringTok}[1]{\textcolor[rgb]{0.13,0.47,0.30}{#1}}
\newcommand{\VariableTok}[1]{\textcolor[rgb]{0.07,0.07,0.07}{#1}}
\newcommand{\VerbatimStringTok}[1]{\textcolor[rgb]{0.13,0.47,0.30}{#1}}
\newcommand{\WarningTok}[1]{\textcolor[rgb]{0.37,0.37,0.37}{\textit{#1}}}

\usepackage{longtable,booktabs,array}
\usepackage{calc} % for calculating minipage widths
% Correct order of tables after \paragraph or \subparagraph
\usepackage{etoolbox}
\makeatletter
\patchcmd\longtable{\par}{\if@noskipsec\mbox{}\fi\par}{}{}
\makeatother
% Allow footnotes in longtable head/foot
\IfFileExists{footnotehyper.sty}{\usepackage{footnotehyper}}{\usepackage{footnote}}
\makesavenoteenv{longtable}
\usepackage{graphicx}
\makeatletter
\newsavebox\pandoc@box
\newcommand*\pandocbounded[1]{% scales image to fit in text height/width
  \sbox\pandoc@box{#1}%
  \Gscale@div\@tempa{\textheight}{\dimexpr\ht\pandoc@box+\dp\pandoc@box\relax}%
  \Gscale@div\@tempb{\linewidth}{\wd\pandoc@box}%
  \ifdim\@tempb\p@<\@tempa\p@\let\@tempa\@tempb\fi% select the smaller of both
  \ifdim\@tempa\p@<\p@\scalebox{\@tempa}{\usebox\pandoc@box}%
  \else\usebox{\pandoc@box}%
  \fi%
}
% Set default figure placement to htbp
\def\fps@figure{htbp}
\makeatother





\setlength{\emergencystretch}{3em} % prevent overfull lines

\providecommand{\tightlist}{%
  \setlength{\itemsep}{0pt}\setlength{\parskip}{0pt}}



 


\makeatletter
\@ifpackageloaded{caption}{}{\usepackage{caption}}
\AtBeginDocument{%
\ifdefined\contentsname
  \renewcommand*\contentsname{Table of contents}
\else
  \newcommand\contentsname{Table of contents}
\fi
\ifdefined\listfigurename
  \renewcommand*\listfigurename{List of Figures}
\else
  \newcommand\listfigurename{List of Figures}
\fi
\ifdefined\listtablename
  \renewcommand*\listtablename{List of Tables}
\else
  \newcommand\listtablename{List of Tables}
\fi
\ifdefined\figurename
  \renewcommand*\figurename{Figure}
\else
  \newcommand\figurename{Figure}
\fi
\ifdefined\tablename
  \renewcommand*\tablename{Table}
\else
  \newcommand\tablename{Table}
\fi
}
\@ifpackageloaded{float}{}{\usepackage{float}}
\floatstyle{ruled}
\@ifundefined{c@chapter}{\newfloat{codelisting}{h}{lop}}{\newfloat{codelisting}{h}{lop}[chapter]}
\floatname{codelisting}{Listing}
\newcommand*\listoflistings{\listof{codelisting}{List of Listings}}
\makeatother
\makeatletter
\makeatother
\makeatletter
\@ifpackageloaded{caption}{}{\usepackage{caption}}
\@ifpackageloaded{subcaption}{}{\usepackage{subcaption}}
\makeatother
\usepackage{bookmark}
\IfFileExists{xurl.sty}{\usepackage{xurl}}{} % add URL line breaks if available
\urlstyle{same}
\hypersetup{
  pdftitle={Case Study 1: Battery Logistics Network Design},
  colorlinks=true,
  linkcolor={blue},
  filecolor={Maroon},
  citecolor={Blue},
  urlcolor={Blue},
  pdfcreator={LaTeX via pandoc}}


\title{Case Study 1: Battery Logistics Network Design}
\usepackage{etoolbox}
\makeatletter
\providecommand{\subtitle}[1]{% add subtitle to \maketitle
  \apptocmd{\@title}{\par {\large #1 \par}}{}{}
}
\makeatother
\subtitle{Team-Based Case Study --- SCM Bachelor}
\author{}
\date{}
\begin{document}
\maketitle


\section{Introduction and Learning
Objectives}\label{introduction-and-learning-objectives}

This case study requires teams of 3--4 students to design a logistics
network for lithium-ion battery distribution across northern Germany and
Benelux. You will apply \textbf{three core methods} from the SCM module:

\begin{enumerate}
\def\labelenumi{\arabic{enumi}.}
\tightlist
\item
  \textbf{Warehouse Location Problem (WLP)}: decide which subset of
  candidate warehouses to open
\item
  \textbf{Transportation Problem (TP)}: plan operational shipments from
  warehouses to plants
\item
  \textbf{Inventory and Safety Stock Analysis}: quantify the cost of
  demand uncertainty
\end{enumerate}

On completion, students will be able to:

\begin{itemize}
\tightlist
\item
  Formulate and solve a WLP using both heuristic and exact methods
\item
  Implement the Haversine formula to derive transport cost estimates
  from geographic data
\item
  Apply Vogel's Approximation Method and compare it to an MILP solution
\item
  Quantify the value of inventory pooling under uncertain demand
\end{itemize}

\textbf{Submission}: Written report (max 12 pages) + R script. Team
presentation: 10 minutes.

\begin{center}\rule{0.5\linewidth}{0.5pt}\end{center}

\section{Background}\label{background}

A third-party logistics provider (3PL) has secured a contract to supply
lithium-ion battery modules to automotive assembly plants across
northern Germany and the Benelux region. Battery modules arrive by sea
container at the port of Hamburg and must be distributed to eleven
plants via regional warehouses.

The 3PL operates six potential warehouse sites across the region. Each
site requires capital investment and annual operating expenditure to
comply with \textbf{dangerous goods regulations} for lithium-ion
batteries (ADR/IMDG hazard class 9 --- UN3480/UN3481). These compliance
costs are reflected in the annual fixed costs below.

The 3PL must decide:

\begin{enumerate}
\def\labelenumi{\arabic{enumi}.}
\tightlist
\item
  \textbf{Which warehouses to open} (strategic, multi-year horizon)
\item
  \textbf{How to allocate plant deliveries to open warehouses}
  (operational, annual plan)
\item
  \textbf{How much safety stock to hold} given uncertain plant demand
  (tactical)
\end{enumerate}

\begin{center}\rule{0.5\linewidth}{0.5pt}\end{center}

\section{Data}\label{data}

\subsection{Automotive Plant Locations and Annual
Demand}\label{automotive-plant-locations-and-annual-demand}

The following eleven assembly plants have placed framework contracts
with the 3PL. Annual battery demand is in \textbf{MWh/year}
(megawatt-hours --- the standard unit for battery storage capacity in
automotive logistics).

\begin{longtable}[]{@{}
  >{\raggedright\arraybackslash}p{(\linewidth - 10\tabcolsep) * \real{0.0476}}
  >{\raggedright\arraybackslash}p{(\linewidth - 10\tabcolsep) * \real{0.1111}}
  >{\raggedright\arraybackslash}p{(\linewidth - 10\tabcolsep) * \real{0.0952}}
  >{\raggedleft\arraybackslash}p{(\linewidth - 10\tabcolsep) * \real{0.2222}}
  >{\raggedleft\arraybackslash}p{(\linewidth - 10\tabcolsep) * \real{0.2540}}
  >{\raggedleft\arraybackslash}p{(\linewidth - 10\tabcolsep) * \real{0.2698}}@{}}
\toprule\noalign{}
\begin{minipage}[b]{\linewidth}\raggedright
\#
\end{minipage} & \begin{minipage}[b]{\linewidth}\raggedright
Plant
\end{minipage} & \begin{minipage}[b]{\linewidth}\raggedright
City
\end{minipage} & \begin{minipage}[b]{\linewidth}\raggedleft
Latitude (°N)
\end{minipage} & \begin{minipage}[b]{\linewidth}\raggedleft
Longitude (°E)
\end{minipage} & \begin{minipage}[b]{\linewidth}\raggedleft
Demand (MWh/yr)
\end{minipage} \\
\midrule\noalign{}
\endhead
\bottomrule\noalign{}
\endlastfoot
1 & VW Wolfsburg & Wolfsburg & 52.432 & 10.787 & 28,000 \\
2 & VW Hannover & Hannover & 52.380 & 9.746 & 22,000 \\
3 & VW Emden & Emden & 53.365 & 7.213 & 18,500 \\
4 & VW Osnabrück & Osnabrück & 52.279 & 8.047 & 12,000 \\
5 & Mercedes-Benz Bremen & Bremen & 53.075 & 8.808 & 25,000 \\
6 & Ford Köln & Köln & 50.950 & 6.900 & 20,000 \\
7 & Opel Rüsselsheim & Rüsselsheim & 49.993 & 8.422 & 15,000 \\
8 & Volvo Cars Gent & Gent (BE) & 51.054 & 3.717 & 30,000 \\
9 & Volvo Trucks Gent & Gent (BE) & 51.040 & 3.730 & 14,000 \\
10 & Toyota Zaventem & Zaventem (BE) & 50.898 & 4.480 & 16,000 \\
11 & VDL Nedcar Born & Born (NL) & 51.030 & 5.867 & 11,000 \\
\end{longtable}

\textbf{Total annual demand: 211,500 MWh/year}

\subsection{Candidate Warehouse Locations and Fixed
Costs}\label{candidate-warehouse-locations-and-fixed-costs}

\begin{longtable}[]{@{}
  >{\raggedright\arraybackslash}p{(\linewidth - 10\tabcolsep) * \real{0.0417}}
  >{\raggedright\arraybackslash}p{(\linewidth - 10\tabcolsep) * \real{0.1389}}
  >{\raggedright\arraybackslash}p{(\linewidth - 10\tabcolsep) * \real{0.0833}}
  >{\raggedleft\arraybackslash}p{(\linewidth - 10\tabcolsep) * \real{0.1944}}
  >{\raggedleft\arraybackslash}p{(\linewidth - 10\tabcolsep) * \real{0.2222}}
  >{\raggedleft\arraybackslash}p{(\linewidth - 10\tabcolsep) * \real{0.3194}}@{}}
\toprule\noalign{}
\begin{minipage}[b]{\linewidth}\raggedright
\#
\end{minipage} & \begin{minipage}[b]{\linewidth}\raggedright
Location
\end{minipage} & \begin{minipage}[b]{\linewidth}\raggedright
City
\end{minipage} & \begin{minipage}[b]{\linewidth}\raggedleft
Latitude (°N)
\end{minipage} & \begin{minipage}[b]{\linewidth}\raggedleft
Longitude (°E)
\end{minipage} & \begin{minipage}[b]{\linewidth}\raggedleft
Annual Fixed Cost (€K)
\end{minipage} \\
\midrule\noalign{}
\endhead
\bottomrule\noalign{}
\endlastfoot
1 & Gütersloh & Gütersloh & 51.906 & 8.383 & 450 \\
2 & Düren & Düren & 50.804 & 6.493 & 480 \\
3 & Dorsten/Marl & Dorsten & 51.660 & 6.965 & 520 \\
4 & Harsewinkel & Harsewinkel & 51.960 & 8.230 & 470 \\
5 & Langenhagen & Langenhagen & 52.438 & 9.737 & 510 \\
6 & Hannover-Süd & Hannover & 52.350 & 9.800 & 460 \\
\end{longtable}

\subsection{Transport Cost Derivation}\label{transport-cost-derivation}

Students derive the transport cost matrix using the \textbf{Haversine
formula}. The rate agreed with the carrier is:

\[c_{ij} = 0.08 \text{ €/km/MWh} \times d_{ij}^{\text{Haversine}} \text{ km}\]

This gives a \textbf{cost per MWh per year} for each warehouse-to-plant
lane. The R code skeleton below computes this matrix --- complete it as
part of your analysis:

\begin{Shaded}
\begin{Highlighting}[]
\NormalTok{haversine }\OtherTok{\textless{}{-}} \ControlFlowTok{function}\NormalTok{(lat1, lon1, lat2, lon2) \{}
\NormalTok{  R    }\OtherTok{\textless{}{-}} \DecValTok{6371}
\NormalTok{  phi1 }\OtherTok{\textless{}{-}}\NormalTok{ lat1 }\SpecialCharTok{*}\NormalTok{ pi }\SpecialCharTok{/} \DecValTok{180}
\NormalTok{  phi2 }\OtherTok{\textless{}{-}}\NormalTok{ lat2 }\SpecialCharTok{*}\NormalTok{ pi }\SpecialCharTok{/} \DecValTok{180}
\NormalTok{  dphi }\OtherTok{\textless{}{-}}\NormalTok{ (lat2 }\SpecialCharTok{{-}}\NormalTok{ lat1) }\SpecialCharTok{*}\NormalTok{ pi }\SpecialCharTok{/} \DecValTok{180}
\NormalTok{  dlam }\OtherTok{\textless{}{-}}\NormalTok{ (lon2 }\SpecialCharTok{{-}}\NormalTok{ lon1) }\SpecialCharTok{*}\NormalTok{ pi }\SpecialCharTok{/} \DecValTok{180}
\NormalTok{  a    }\OtherTok{\textless{}{-}} \FunctionTok{sin}\NormalTok{(dphi}\SpecialCharTok{/}\DecValTok{2}\NormalTok{)}\SpecialCharTok{\^{}}\DecValTok{2} \SpecialCharTok{+} \FunctionTok{cos}\NormalTok{(phi1)}\SpecialCharTok{*}\FunctionTok{cos}\NormalTok{(phi2)}\SpecialCharTok{*}\FunctionTok{sin}\NormalTok{(dlam}\SpecialCharTok{/}\DecValTok{2}\NormalTok{)}\SpecialCharTok{\^{}}\DecValTok{2}
\NormalTok{  R }\SpecialCharTok{*} \DecValTok{2} \SpecialCharTok{*} \FunctionTok{atan2}\NormalTok{(}\FunctionTok{sqrt}\NormalTok{(a), }\FunctionTok{sqrt}\NormalTok{(}\DecValTok{1}\SpecialCharTok{{-}}\NormalTok{a))}
\NormalTok{\}}

\CommentTok{\# Warehouses: 6 rows; Plants: 11 columns}
\CommentTok{\# Build the 6×11 transport cost matrix (€/MWh)}
\CommentTok{\# cost\_matrix[i,j] = 0.08 * haversine(wh\_lat[i], wh\_lon[i], plant\_lat[j], plant\_lon[j])}
\end{Highlighting}
\end{Shaded}

\begin{center}\rule{0.5\linewidth}{0.5pt}\end{center}

\section{Part A: Strategic Network Design (Warehouse Location
Problem)}\label{part-a-strategic-network-design-warehouse-location-problem}

\subsection{A.1 Problem Formulation}\label{a.1-problem-formulation}

Formulate the WLP for this case as a binary integer programme. Define:

\begin{itemize}
\tightlist
\item
  Decision variables: \(y_i \in \{0,1\}\) for \(i = 1, \ldots, 6\) (1 =
  open warehouse \(i\))
\item
  Decision variables: \(x_{ij} \in [0,1]\) for all \(i, j\) (fraction of
  plant \(j\)'s demand served from warehouse \(i\))
\item
  Objective function: minimise total annual cost (fixed + transport)
\item
  Constraints: (i) each plant is fully served, (ii) plants can only be
  served from open warehouses
\end{itemize}

Write out the complete formulation with symbols, then translate to R.

\subsection{A.2 Transport Cost Matrix}\label{a.2-transport-cost-matrix}

Using the R skeleton above and the coordinate data provided:

\begin{enumerate}
\def\labelenumi{\arabic{enumi}.}
\tightlist
\item
  Compute all 66 Haversine distances (6 warehouses × 11 plants)
\item
  Multiply by the transport rate €0.08/km/MWh to get the per-unit-demand
  annual cost
\item
  Present the resulting cost matrix as a formatted table in your report
\end{enumerate}

\textbf{Discussion question:} What simplifications does this
distance-based cost model make? How would you refine it for a real
project?

\subsection{A.3 Add Heuristic}\label{a.3-add-heuristic}

Apply the \textbf{Add Heuristic} to the transport cost matrix derived in
A.2:

\begin{enumerate}
\def\labelenumi{\arabic{enumi}.}
\tightlist
\item
  Start with no warehouses open
\item
  At each step, compute the net saving from opening each remaining
  candidate warehouse
\item
  Open the warehouse with the highest net saving (if positive)
\item
  Repeat until no further improvement is possible
\end{enumerate}

\textbf{Present:} the savings table for each iteration, the final set of
open warehouses, and the total cost breakdown (fixed + transport).

\subsection{A.4 Exact MILP Solution}\label{a.4-exact-milp-solution}

Solve the WLP exactly using \texttt{ompr} or \texttt{lpSolve} in R:

\begin{Shaded}
\begin{Highlighting}[]
\CommentTok{\# Sketch (complete this):}
\FunctionTok{library}\NormalTok{(ompr); }\FunctionTok{library}\NormalTok{(ompr.roi); }\FunctionTok{library}\NormalTok{(ROI.plugin.glpk)}
\NormalTok{n\_wh }\OtherTok{\textless{}{-}} \DecValTok{6}\NormalTok{; n\_pl }\OtherTok{\textless{}{-}} \DecValTok{11}

\NormalTok{model }\OtherTok{\textless{}{-}} \FunctionTok{MIPModel}\NormalTok{() }\SpecialCharTok{|\textgreater{}}
  \FunctionTok{add\_variable}\NormalTok{(y[i], }\AttributeTok{i =} \DecValTok{1}\SpecialCharTok{:}\NormalTok{n\_wh, }\AttributeTok{type =} \StringTok{"binary"}\NormalTok{) }\SpecialCharTok{|\textgreater{}}
  \FunctionTok{add\_variable}\NormalTok{(x[i,j], }\AttributeTok{i =} \DecValTok{1}\SpecialCharTok{:}\NormalTok{n\_wh, }\AttributeTok{j =} \DecValTok{1}\SpecialCharTok{:}\NormalTok{n\_pl,}
               \AttributeTok{type =} \StringTok{"continuous"}\NormalTok{, }\AttributeTok{lb =} \DecValTok{0}\NormalTok{, }\AttributeTok{ub =} \DecValTok{1}\NormalTok{) }\SpecialCharTok{|\textgreater{}}
  \FunctionTok{set\_objective}\NormalTok{(...) }\SpecialCharTok{|\textgreater{}}
  \FunctionTok{add\_constraint}\NormalTok{(...) }\SpecialCharTok{|\textgreater{}}
  \FunctionTok{add\_constraint}\NormalTok{(...)}

\NormalTok{result }\OtherTok{\textless{}{-}} \FunctionTok{solve\_model}\NormalTok{(model, }\FunctionTok{with\_ROI}\NormalTok{(}\AttributeTok{solver =} \StringTok{"glpk"}\NormalTok{))}
\end{Highlighting}
\end{Shaded}

Report the optimal set of open warehouses and the total cost.

\subsection{A.5 Comparison: Heuristic
vs.~Optimal}\label{a.5-comparison-heuristic-vs.-optimal}

\begin{longtable}[]{@{}lcc@{}}
\toprule\noalign{}
Criterion & Add Heuristic & MILP Optimal \\
\midrule\noalign{}
\endhead
\bottomrule\noalign{}
\endlastfoot
Open warehouses & ? & ? \\
Fixed cost (€K/yr) & ? & ? \\
Transport cost (€K/yr) & ? & ? \\
\textbf{Total cost (€K/yr)} & \textbf{?} & \textbf{?} \\
Optimality gap & --- & 0\% \\
\end{longtable}

Answer: \textbf{How close is the heuristic?} In what situations might
the gap be larger?

\begin{center}\rule{0.5\linewidth}{0.5pt}\end{center}

\section{Part B: Operational Transport
Planning}\label{part-b-operational-transport-planning}

\subsection{Scenario}\label{scenario}

The 3PL has reached a strategic decision (based on regulatory and
operational considerations) to open the following three warehouses:

\begin{longtable}[]{@{}llc@{}}
\toprule\noalign{}
Warehouse & Location & Capacity (MWh/yr) \\
\midrule\noalign{}
\endhead
\bottomrule\noalign{}
\endlastfoot
Düren & Düren & 150,000 \\
Harsewinkel & Harsewinkel & 50,000 \\
Langenhagen & Langenhagen & 100,000 \\
\end{longtable}

\textbf{Total warehouse capacity: 300,000 MWh/yr \textgreater{} 211,500
MWh/yr demand} (system is feasible)

For operational planning, the annual demand figures have been updated
following a new forecast round:

\begin{longtable}[]{@{}lc@{}}
\toprule\noalign{}
Plant & Updated Demand (MWh/yr) \\
\midrule\noalign{}
\endhead
\bottomrule\noalign{}
\endlastfoot
VW Wolfsburg & 29,500 \\
VW Hannover & 21,000 \\
VW Emden & 19,000 \\
VW Osnabrück & 13,000 \\
Mercedes-Benz Bremen & 26,000 \\
Ford Köln & 21,500 \\
Opel Rüsselsheim & 14,500 \\
Volvo Cars Gent & 31,000 \\
Volvo Trucks Gent & 14,500 \\
Toyota Zaventem & 16,500 \\
VDL Nedcar Born & 12,500 \\
\textbf{Total} & \textbf{219,000} \\
\end{longtable}

\subsection{B.1 Transportation Problem
Formulation}\label{b.1-transportation-problem-formulation}

Write out the balanced transportation problem. Because total capacity
(300,000) exceeds total demand (219,000), you must add a dummy plant
with demand = 81,000 MWh and zero transport cost to all warehouses.

The cost matrix for this sub-problem should be derived from the
Haversine distances computed in Part A.

\subsection{B.2 Vogel's Approximation
Method}\label{b.2-vogels-approximation-method}

Apply VAM to the 3×12 (including dummy) transportation problem:

\begin{enumerate}
\def\labelenumi{\arabic{enumi}.}
\tightlist
\item
  Compute row and column penalties at each step
\item
  Allocate to the cell with minimum cost in the highest-penalty row or
  column
\item
  Eliminate exhausted supply/demand and repeat
\end{enumerate}

Show all intermediate steps in tabular form. Report the total
transportation cost of your VAM solution.

\subsection{B.3 Solve with lpSolve}\label{b.3-solve-with-lpsolve}

\begin{Shaded}
\begin{Highlighting}[]
\FunctionTok{library}\NormalTok{(lpSolve)}
\CommentTok{\# Use lp.transport() — see Exercise Block 3 for reference}
\NormalTok{result }\OtherTok{\textless{}{-}} \FunctionTok{lp.transport}\NormalTok{(}
  \AttributeTok{cost.mat  =}\NormalTok{ cost\_sub,       }\CommentTok{\# 3 x 12 (incl. dummy)}
  \AttributeTok{direction =} \StringTok{"min"}\NormalTok{,}
  \AttributeTok{row.signs =} \FunctionTok{rep}\NormalTok{(}\StringTok{"\textless{}="}\NormalTok{, }\DecValTok{3}\NormalTok{),}
  \AttributeTok{row.rhs   =}\NormalTok{ capacity\_vec,}
  \AttributeTok{col.signs =} \FunctionTok{rep}\NormalTok{(}\StringTok{"\textgreater{}="}\NormalTok{, }\DecValTok{12}\NormalTok{),}
  \AttributeTok{col.rhs   =}\NormalTok{ demand\_vec\_bal}
\NormalTok{)}
\end{Highlighting}
\end{Shaded}

Report the optimal allocation matrix and total cost.

\subsection{B.4 Discussion}\label{b.4-discussion}

\begin{enumerate}
\def\labelenumi{\arabic{enumi}.}
\tightlist
\item
  Compare VAM and LP-optimal solutions. What is the gap?
\item
  Is the chosen warehouse configuration (Düren/Harsewinkel/Langenhagen)
  cost-efficient? Would a different trio from Part A yield lower
  transport costs?
\item
  Langenhagen has the highest fixed cost among the three. Under what
  conditions would it be worth replacing with Hannover-Süd (slightly
  lower fixed cost, similar location)?
\end{enumerate}

\begin{center}\rule{0.5\linewidth}{0.5pt}\end{center}

\section{Part C: Uncertainty and
Risk}\label{part-c-uncertainty-and-risk}

\subsection{C.1 Demand Uncertainty
Model}\label{c.1-demand-uncertainty-model}

Battery demand from automotive plants is \textbf{uncertain}: programme
ramp-ups, model changeovers, and supply chain disruptions all introduce
variability. Based on historical data from comparable programmes, the
3PL's planning team models each plant's annual demand as
\textbf{normally distributed}:

\[D_j \sim \mathcal{N}(\mu_j, \, \sigma_j^2) \quad \text{with} \quad \sigma_j = 0.15 \cdot \mu_j\]

This corresponds to a \textbf{coefficient of variation of 15\%} --- a
moderate level of forecast uncertainty typical of 12-month rolling
forecasts.

\textbf{Task:} For each of the three open warehouses, compute the
\textbf{aggregate demand} from its assigned plants (using the LP
allocation from B.3). What is the standard deviation of this aggregate
demand, assuming plant demands are \textbf{independent}?

\[\sigma_{\text{WH}_k} = \sqrt{\sum_{j \in \mathcal{J}_k} \sigma_j^2}\]

\subsection{C.2 Safety Stock at Each
Warehouse}\label{c.2-safety-stock-at-each-warehouse}

For a \textbf{95\% type-I service level} (probability of no stockout in
a planning period), the required safety stock at warehouse \(k\) is:

\[SS_k = z_{0.95} \cdot \sigma_{\text{WH}_k} \cdot \sqrt{L}\]

where \(L\) is the replenishment lead time (assume \(L = 2\) weeks =
2/52 year) and \(z_{0.95} = 1.645\).

\begin{enumerate}
\def\labelenumi{\arabic{enumi}.}
\tightlist
\item
  Compute \(SS_k\) for each of the three warehouses
\item
  Convert safety stock from MWh to monetary value (assume battery cost
  €80/kWh = €80,000/MWh)
\item
  Sum the total safety stock investment across all three warehouses
\end{enumerate}

\subsection{C.3 Centralised Inventory
Alternative}\label{c.3-centralised-inventory-alternative}

Consider the hypothetical alternative of supplying all 11 plants from a
\textbf{single central warehouse} (e.g.~located at the geometric
centroid of all plants).

Under full centralisation, the aggregate standard deviation becomes:

\[\sigma_{\text{central}} = \sqrt{\sum_{j=1}^{11} \sigma_j^2}\]

(independent demands). Compare this to the sum
\(\sum_k \sigma_{\text{WH}_k}\) from C.2.

\textbf{Tasks:}

\begin{enumerate}
\def\labelenumi{\arabic{enumi}.}
\tightlist
\item
  Compute the safety stock under centralisation (same 95\% SL, same lead
  time)
\item
  Compare total safety stock investment: decentralised network
  vs.~centralised alternative
\item
  The \textbf{pooling factor} is \(\sqrt{n_{\text{locations}}}\) ---
  verify this holds approximately here
\end{enumerate}

\textbf{Discussion questions:}

\begin{itemize}
\tightlist
\item
  What trade-offs does centralisation imply for transport costs, lead
  times, and service levels?
\item
  In the context of battery logistics (high unit value, hazardous
  goods), which argument weighs more heavily: pooling savings or
  proximity to customers?
\item
  How would the analysis change if plant demands were
  \textbf{correlated} (e.g.~due to a common automotive market cycle)?
\end{itemize}

\begin{center}\rule{0.5\linewidth}{0.5pt}\end{center}

\section{Grading Rubric}\label{grading-rubric}

\begin{longtable}[]{@{}
  >{\raggedright\arraybackslash}p{(\linewidth - 4\tabcolsep) * \real{0.2075}}
  >{\centering\arraybackslash}p{(\linewidth - 4\tabcolsep) * \real{0.1509}}
  >{\raggedright\arraybackslash}p{(\linewidth - 4\tabcolsep) * \real{0.6415}}@{}}
\toprule\noalign{}
\begin{minipage}[b]{\linewidth}\raggedright
Criterion
\end{minipage} & \begin{minipage}[b]{\linewidth}\centering
Weight
\end{minipage} & \begin{minipage}[b]{\linewidth}\raggedright
Indicators of Strong Performance
\end{minipage} \\
\midrule\noalign{}
\endhead
\bottomrule\noalign{}
\endlastfoot
Problem formulation correctness & 25\% & Correct MILP formulation;
balanced TP; proper use of Haversine; no logical errors in model
setup \\
Solution methodology & 25\% & VAM steps shown in full; lpSolve used
correctly; heuristic and MILP compared; safety stock formula applied
correctly \\
Interpretation and management recommendations & 30\% & Logical
conclusions drawn from results; trade-offs discussed; sensitivity to
assumptions addressed; practical relevance of findings \\
Presentation quality & 20\% & Clear slides, results visualised; team
presents confidently; answers audience questions; report is
well-structured and within page limit \\
\end{longtable}

\textbf{Pass threshold}: minimum 50\% total, with at least 40\% in each
individual criterion.

\begin{center}\rule{0.5\linewidth}{0.5pt}\end{center}

\section{Appendix A: Useful R
Packages}\label{appendix-a-useful-r-packages}

\begin{Shaded}
\begin{Highlighting}[]
\CommentTok{\# Geographic distances}
\CommentTok{\# (no additional package needed — implement Haversine manually)}

\CommentTok{\# Transportation problem}
\FunctionTok{install.packages}\NormalTok{(}\StringTok{"lpSolve"}\NormalTok{)}
\FunctionTok{library}\NormalTok{(lpSolve)}

\CommentTok{\# MILP / WLP}
\FunctionTok{install.packages}\NormalTok{(}\FunctionTok{c}\NormalTok{(}\StringTok{"ompr"}\NormalTok{, }\StringTok{"ompr.roi"}\NormalTok{, }\StringTok{"ROI"}\NormalTok{, }\StringTok{"ROI.plugin.glpk"}\NormalTok{))}
\FunctionTok{library}\NormalTok{(ompr); }\FunctionTok{library}\NormalTok{(ompr.roi); }\FunctionTok{library}\NormalTok{(ROI.plugin.glpk)}

\CommentTok{\# Visualisation}
\FunctionTok{library}\NormalTok{(ggplot2)}
\FunctionTok{install.packages}\NormalTok{(}\StringTok{"sf"}\NormalTok{)       }\CommentTok{\# optional: for proper geographic maps}
\FunctionTok{install.packages}\NormalTok{(}\StringTok{"ggrepel"}\NormalTok{)  }\CommentTok{\# non{-}overlapping text labels}
\end{Highlighting}
\end{Shaded}

\begin{center}\rule{0.5\linewidth}{0.5pt}\end{center}

\section{Appendix B: Report Structure
Guidance}\label{appendix-b-report-structure-guidance}

Your written report should follow this structure (max 12 pages,
excluding appendix):

\begin{enumerate}
\def\labelenumi{\arabic{enumi}.}
\tightlist
\item
  \textbf{Executive Summary} (max 1 page): key recommendation and
  headline cost figures
\item
  \textbf{Part A: Network Design} --- formulation, cost matrix,
  heuristic vs.~MILP, recommendation
\item
  \textbf{Part B: Transport Planning} --- TP setup, VAM steps (summary
  table), LP solution, discussion
\item
  \textbf{Part C: Risk and Inventory} --- safety stock calculation,
  pooling comparison, managerial insight
\item
  \textbf{Conclusion}: Which network do you recommend? What are the key
  uncertainties?
\item
  \textbf{Appendix}: Full R code (not counted in 12 pages)
\end{enumerate}

\textbf{Note:} All figures must have a caption and be referenced in the
text. All tables must be formatted (not R console output). Quantitative
results must be explicitly stated in the text, not just shown in a
table.




\end{document}
